% Copyright (C) 2001-2010 Roberto Bagnara <bagnara@cs.unipr.it>
% Copyright (C) 2010-2020 BUGSENG srl (http://bugseng.com)
%
% This document describes the Parma Polyhedra Library (PPL).
%
% Permission is granted to copy, distribute and/or modify this document
% under the terms of the GNU Free Documentation License, Version 1.2
% or any later version published by the Free Software Foundation;
% with no Invariant Sections, no Front-Cover Texts, and no Back-Cover Texts.
% The license is included, in various formats, in the `doc' subdirectory
% of each distribution of the PPL in files named `fdl.*'.
%
% The PPL is free software; you can redistribute it and/or modify it
% under the terms of the GNU General Public License as published by the
% Free Software Foundation; either version 3 of the License, or (at your
% option) any later version.  The license is included, in various
% formats, in the `doc' subdirectory of each distribution of the PPL in
% files named `gpl.*'.
%
% The PPL is distributed in the hope that it will be useful, but WITHOUT
% ANY WARRANTY; without even the implied warranty of MERCHANTABILITY or
% FITNESS FOR A PARTICULAR PURPOSE.  See the GNU General Public License
% for more details.
%
% If you have not received a copy of one or both the above mentioned
% licenses along with the PPL, write to the Free Software Foundation,
% Inc., 51 Franklin Street, Fifth Floor, Boston, MA 02111-1307, USA.
%
% For the most up-to-date information see the Parma Polyhedra Library
% site: http://bugseng.com/products/ppl/ .

\documentclass[a4paper,twoside]{article}
\usepackage{a4wide}
\usepackage{makeidx}
\usepackage{fancyhdr}
\usepackage{graphicx}
\usepackage{multicol}
\usepackage{float}
\usepackage{textcomp}
\usepackage{alltt}
\usepackage{times}
\usepackage{tocloft}
\usepackage{ifpdf}
\usepackage[table]{xcolor}
\ifpdf
\usepackage[pdftex,
            pagebackref=true,
            colorlinks=true,
            linkcolor=blue,
            unicode
           ]{hyperref}
\else
\usepackage[ps2pdf,
            pagebackref=true,
            colorlinks=true,
            linkcolor=blue,
            unicode
           ]{hyperref}
\usepackage{pspicture}
\fi
\usepackage[utf8]{inputenc}
\usepackage{doxygen}
\newcommand{\+}{\discretionary{\mbox{\scriptsize$\hookleftarrow$}}{}{}}
\usepackage{ppl}
\makeindex
\setcounter{tocdepth}{2}

\renewcommand{\footrulewidth}{0.4pt}
\fancyfoot{}
\fancyfoot[LE,RO]{\includegraphics[height=0.4cm]{ppl_logo_no_ppl.pdf}}
\fancyfoot[RE,LO]{\small The PPL User's Manual (version $projectnumber). See \url{http://bugseng.com} for more information.}

\setlength{\headheight}{24pt}

% Custom commands
\newcommand{\clearemptydoublepage}{%
  \newpage{\pagestyle{empty}\cleardoublepage}%
}

\begin{document}
\title{
\includegraphics[height=9cm]{ppl_logo.pdf} \\
\ \\
The Parma Polyhedra Library \\
User's Manual\thanks{This work is based on previous work also by
Elisa Ricci, Sara Bonini, Andrea Pescetti, Angela Stazzone, Tatiana Zolo.
This work has been partly supported by:
  University of Parma's FIL scientific research project (ex 60\%)
    ``Pure and Applied Mathematics'';
  MURST project
    ``Automatic Program Certification by Abstract Interpretation'';
  MURST project
    ``Abstract Interpretation, Type Systems and Control-Flow Analysis'';
  MURST project
    ``Automatic Aggregate- and Number-Reasoning for Computing: from
      Decision Algorithms to Constraint Programming with Multisets,
      Sets, and Maps'';
  MURST project
    ``Constraint Based Verification of Reactive Systems'';
  MURST project
    ``Abstract Interpretation: Design and Applications'';
  EPSRC project
    ``Numerical Domains for Software Analysis'';
  EPSRC project
    ``Geometric Abstractions for Scalable Program Analyzers''.
  } \\
(version $projectnumber)
}
\author{
Roberto Bagnara\thanks{bagnara@cs.unipr.it,
  Department of Mathematics, University of Parma, Italy,
  and BUGSENG srl.} \\
Patricia M. Hill\thanks{patricia.hill@bugseng.com,
  BUGSENG srl.} \\
Enea Zaffanella\thanks{zaffanella@cs.unipr.it,
  Department of Mathematics, University of Parma, Italy,
  and BUGSENG srl.} \\
Abramo Bagnara\thanks{abramo.bagnara@bugseng.com,
  BUGSENG srl.} \\
}
{\sffamily\maketitle}
\thispagestyle{empty}
\clearpage{\pagestyle{empty}\cleardoublepage}

\newpage
Copyright \copyright\ 2001--2010 Roberto Bagnara (bagnara@cs.unipr.it) \\
Copyright \copyright\ 2010--2020 BUGSENG srl (http://bugseng.com)


This document describes the Parma Polyhedra Library (PPL).

Permission is granted to copy, distribute and/or modify this document
under the terms of the GNU Free Documentation License, Version 1.2
or any later version published by the
\href{http://www.fsf.org}{Free Software Foundation};
with no Invariant Sections, no Front-Cover Texts, and no Back-Cover Texts.
A copy of the license is included in the section entitled
``\hyperlink{GFDL_GFDL}{GNU Free Documentation License}''.

The PPL is free software; you can redistribute it and/or modify it
under the terms of the GNU General Public License as published by the
\href{http://www.fsf.org}{Free Software Foundation}; either version 3
of the License, or (at your option) any later version.
A copy of the license is included in the section entitled
``\hyperlink{GPL_GPL}{GNU GENERAL PUBLIC LICENSE}''.

The PPL is distributed in the hope that it will be useful, but WITHOUT
ANY WARRANTY; without even the implied warranty of MERCHANTABILITY or
FITNESS FOR A PARTICULAR PURPOSE.  See the GNU General Public License
for more details.

If you have not received a copy of one or both the above mentioned
licenses along with the PPL, write to the Free Software Foundation,
Inc., 51 Franklin Street, Fifth Floor, Boston, MA 02111-1307, USA.

For the most up-to-date information see the Parma Polyhedra Library
site:
\begin{center}
\href{http://bugseng.com/products/ppl/}{\tt http://bugseng.com/products/ppl/}
\end{center}
\vfill
\begin{center}
\includegraphics[width=5cm]{bugseng_logo.pdf}
\end{center}
\vfill
\clearpage{\pagestyle{empty}\cleardoublepage}

% The following command is needed to have the right footer on
% the first page of the TOC.
\addtocontents{toc}{\protect\thispagestyle{fancy}}

% Make sure there is enough space in the TOC for long subsection numbers.
\addtolength{\cftsubsecnumwidth}{1em}

\pagenumbering{roman}
\tableofcontents
\clearpage{\pagestyle{empty}\cleardoublepage}

\pagenumbering{arabic}
